\chapter{Objetivos}

\section{Objetivo general}

Diseñar y desarrollar un prototipo funcional de software que analice resonancias magnéticas lumbares en los planos sagital y axial para segmentar estructuras anatómicas relevantes, obtener un conjunto acotado de mediciones cuantitativas y generar una salida visual y estructurada editable (incluido un preinforme automático revisable) que asista al profesional sin reemplazar su criterio clínico.

\section{Objetivos específicos}

Para alcanzar el objetivo general, se definen los siguientes objetivos específicos:

\begin{enumerate}
  \item Relevar antecedentes académicos, datasets públicos, herramientas abiertas y soluciones comerciales relacionadas con el análisis asistido de RM lumbar.
  \item Identificar la brecha funcional que justifica un prototipo académico orientado a segmentación, mediciones, visualización y salida editable.
  \item Definir el alcance del MVP, incluyendo estructuras anatómicas, funcionalidades, validación prevista y exclusiones clínicas, técnicas y regulatorias.
  \item Analizar necesidades de potenciales usuarios mediante actividades de user research vinculadas con la revisión de RM lumbar.
  \item Diseñar e implementar un pipeline de carga, normalización, preprocesamiento, segmentación, postprocesamiento y visualización.
  \item Adaptar o implementar un modelo de segmentación para vértebras, discos intervertebrales y canal espinal en el plano sagital, y un módulo complementario de segmentación axial (disco, elemento posterior, saco tecal y área entre elementos anterior y posterior).
  \item Desarrollar mediciones geométricas simples derivadas de las máscaras generadas.
  \item Diseñar una interfaz, una salida estructurada editable y un preinforme automático revisable que permitan revisar resultados sin constituir un informe clínico automático.
  \item Evaluar técnicamente el prototipo mediante métricas de segmentación y comparación contra anotaciones de referencia.
  \item Complementar la validación técnica con revisión cualitativa profesional y documentar decisiones, limitaciones, resultados y trabajos futuros.
\end{enumerate}
